\section{Estudio de Factibilidad y Gestión de Riesgos}

El presente estudio analiza la viabilidad del proyecto desde cuatro dimensiones
complementarias: técnica, operativa, económica y legal. El objetivo es demostrar,
con evidencia cuantitativa y referencias documentadas, que el proyecto puede
construirse con los recursos disponibles, ser adoptado por sus usuarios finales,
mantenerse operativamente con costos controlados y operar dentro del marco
normativo vigente en México.

La metodología aplicada combina: (1) revisión de documentación oficial de las
plataformas tecnológicas seleccionadas; (2) análisis de estadísticas de adopción
digital publicadas por el INEGI y organismos de la industria; (3) modelado de
costos basado en tarifas públicas vigentes; y (4) revisión del marco normativo
mexicano aplicable al manejo de datos personales, registros internos de liquidación,
servicios digitales y software \cite{nextjs, supabase_pricing, prisma_docs, state_of_js_2023, stackoverflow_2023_results}.

% -----------------------------------------------------------
\FloatBarrier
\subsection{Factibilidad Técnica}
% -----------------------------------------------------------

La factibilidad técnica evalúa la disponibilidad y madurez de las herramientas,
lenguajes, frameworks e infraestructura necesarios para el desarrollo e implantación
del sistema.

\subsubsection{Ecosistema de Desarrollo: Madurez y Adopción}

La selección del stack tecnológico se fundamenta en el criterio de madurez
comprobada, adopción masiva en la industria y soporte activo de la comunidad.
El uso de tecnologías ampliamente adoptadas reduce el riesgo técnico, amplía el
pool de desarrolladores disponibles y garantiza continuidad del soporte a largo
plazo. La Tabla~\ref{tab:stack} resume las tecnologías principales con sus métricas
de adopción.

\begin{table}[H]
\caption{Stack Tecnológico: Madurez y Adopción}
\label{tab:stack}
\centering
\small
\setlength{\tabcolsep}{3pt}
\begin{tabular}{>{\raggedright\arraybackslash}p{2.3cm}
                >{\raggedright\arraybackslash}p{3cm}
                >{\raggedright\arraybackslash}p{1.5cm}
                >{\raggedright\arraybackslash}p{6cm}
                >{\raggedright\arraybackslash}p{2.5cm}}
\hline
\thead{Tecnología} & \thead{Categoría} & \thead{Versión} & \thead{Adopción} & \thead{Licencia} \\
\hline
Next.js 16 & Framework frontend + API & 16.x & 44\% de desarrolladores lo usan activamente \cite{state_of_js_2023} & MIT \\ \rowcolor{altrow}
React 19 & Biblioteca UI & 19.x & 84\% de los desarrolladores web la usan \cite{state_of_js_2023} & MIT \\
TypeScript 5 & JS tipado & 5.x & 38.87\% de desarrolladores la usan \cite{stackoverflow_2023_results} & Apache 2.0 \\ \rowcolor{altrow}
Supabase & BaaS: PostgreSQL + Auth + Realtime & GA 2023 & +65,000 estrellas en GitHub; 200,000+ proyectos activos \cite{supabase_pricing} & Apache 2.0 \\
PostgreSQL 16 & SGBD relacional & 16.x & 45.55\% de popularidad \cite{stackoverflow_2023_results} & PostgreSQL License \\ \rowcolor{altrow}
Prisma ORM & ORM para Node.js/TS & Versión estable vigente & Herramienta ampliamente utilizada en ecosistemas TypeScript y Next.js \cite{prisma_docs} & Apache 2.0 \\
Tailwind CSS 4 & Framework CSS & 4.x & Framework CSS líder por su enfoque utility-first \cite{state_of_js_2023} & MIT \\ \rowcolor{altrow}
shadcn/ui 4 & Biblioteca de componentes & 4.x & Biblioteca ampliamente utilizada para construir componentes accesibles y reutilizables & MIT \\
Vercel & Hosting serverless & PaaS & Plataforma l\'ider para despliegue Next.js \cite{vercel_pricing, nextjs} & Propietario \\ \rowcolor{altrow}
Node.js 20 LTS & Runtime JS servidor & 20.x LTS & framework m\'as usado (42.65\%) \cite{stackoverflow_2023_results} & MIT \\
\hline
\end{tabular}
\end{table}

\noindent\footnotesize{\textit{Nota:} Las métricas de adopción se usan como referencia general del ecosistema tecnológico, no como indicador específico de la versión exacta seleccionada.}
\normalsize

\subsubsection{Análisis de Componentes Críticos}

\noindent\textbf{Vercel como plataforma de hosting.}
Vercel es la plataforma de referencia para despliegue de aplicaciones Next.js,
siendo el creador y principal mantenedor del framework. Ofrece escalado automático,
red CDN global en más de 100 edge locations, deploy automático por rama de Git y
ambientes de Preview y Producción nativos. El plan Pro ofrece límites superiores
al plan gratuito y resulta adecuado para endpoints breves de la aplicación, mientras
que los procesos largos se delegan al pipeline batch ejecutado en el VPS
\cite{vercel_pricing, nextjs}.

\noindent\textbf{Supabase como backend-as-a-service.}
Supabase provee una instancia PostgreSQL administrada junto con servicios de
autenticación, Realtime (basado en LISTEN/NOTIFY de PostgreSQL), almacenamiento
de archivos y una API REST auto-generada. Su elección está fundamentada en que
PostgreSQL es un sistema de base de datos relacional ampliamente adoptado.
El plan Pro incluye almacenamiento administrado, respaldos automáticos y
mecanismos de pooling de conexiones mediante PgBouncer, lo que reduce el riesgo de
agotamiento de conexiones en entornos serverless \cite{supabase_pricing}.

\noindent\textbf{Prisma ORM como capa de acceso a datos.}
Prisma provee un cliente de base de datos tipado generado automáticamente a partir
del schema definido en su DSL. Esto garantiza seguridad de tipos en tiempo de
compilación, previene SQL Injection por diseño al utilizar consultas parametrizadas
y facilita las migraciones de schema con control de versiones. En entornos
serverless, Prisma Accelerate resuelve el problema de agotamiento de conexiones
al actuar como proxy entre las funciones de Vercel y la instancia PostgreSQL
de Supabase \cite{prisma_docs, prisma_accelerate}. Las versiones recientes de Prisma han mejorado el soporte para
entornos serverless y la generación del cliente.

La Tabla~\ref{tab:hw_usuarios} detalla los requerimientos mínimos de hardware
para cada tipo de usuario final.

\begin{table}[H]
\caption{Requerimientos de Hardware por Tipo de Usuario}
\label{tab:hw_usuarios}
\centering
\small
\setlength{\tabcolsep}{3pt}
\begin{tabular}{>{\raggedright\arraybackslash}p{2.3cm}
                >{\raggedright\arraybackslash}p{2.5cm}
                >{\raggedright\arraybackslash}p{5.5cm}
                >{\raggedright\arraybackslash}p{5cm}}
\hline
\thead{Tipo de usuario} & \thead{Dispositivo} & \thead{Requerimiento mínimo} & \thead{Observación} \\
\hline
Comensal / Anfitrión & Smartphone personal & Chrome $\geq$90, Firefox $\geq$88 o Safari $\geq$14. Cámara para QR. & El 87.3\% de mexicanos de 18--44 años usa smartphone con internet \cite{inegi_endutih_2022}. \\ \rowcolor{altrow}
Mesero & Tablet del restaurante & Chrome $\geq$90, Firefox $\geq$88 o Safari $\geq$14. Pantalla $\geq$8 pulgadas recomendada. & Tablet de gama media ($\sim{\$4,500 MXN}$). \\
Chef & Tablet del restaurante & Pantalla $\geq$8 pulgadas recomendada para tablero de órdenes. & Tablet o monitor con PC básica. \\ \rowcolor{altrow}
Gerente / Admin & PC, laptop o tablet & Navegador moderno. Pantalla $\geq$11 pulgadas recomendada para Dashboard. & Equipo existente en la mayoría de negocios. \\
Servidor analítico (VPS) & VPS en la nube & 2 vCPU, 8~GB RAM, 100~GB NVMe, 8~TB de ancho de banda. Ubuntu 22.04~LTS. & Hostinger KVM 2: Promocional \$8.99~USD/mes; renovaci\'on \$14.99~USD/mes. \\
\hline
\end{tabular}
\end{table}

\subsubsection{Disponibilidad de Conocimientos Técnicos}

El stack seleccionado (TypeScript + Next.js + PostgreSQL + React) es el más
demandado en el mercado de desarrollo web en México y Latinoamérica. TypeScript
ha sido el lenguaje más deseado por desarrolladores tres años consecutivos
(2021--2023) \cite{state_of_js_2023}; Next.js ocupa el primer lugar entre frameworks full-stack
deseados \cite{state_of_js_2023}; y PostgreSQL es la base de datos relacional más popular globalmente \cite{state_of_js_2023}.
Plataformas como LinkedIn, OCC Mundial y Computrabajo registran cientos de
ofertas activas con este perfil mensualmente.

\subsubsection{Restricciones Técnicas y Mitigación}

La Tabla~\ref{tab:restricciones_tec} resume las restricciones técnicas identificadas
y las soluciones disponibles para cada una.

\begin{table}[H]
\caption{Restricciones Técnicas Identificadas y Mitigación}
\label{tab:restricciones_tec}
\centering
\small
\setlength{\tabcolsep}{3pt}
\begin{tabular}{>{\raggedright\arraybackslash}p{4cm}
                >{\raggedright\arraybackslash}p{4cm}
                >{\raggedright\arraybackslash}p{7.5cm}}
\hline
\thead{Restricción técnica} & \thead{Impacto potencial} & \thead{Solución técnica disponible} \\
\hline
Límite de conexiones PostgreSQL en funciones serverless &
  Pool de conexiones agotado en horas pico &
  Prisma Accelerate o PgBouncer (integrado en Supabase Pro). Solución probada y documentada \cite{prisma_docs, prisma_accelerate}. \\ \rowcolor{altrow}
Límites de ejecución en funciones serverless &
  Operaciones largas o procesos analíticos podrían exceder el tiempo permitido
  por la plataforma &
  Diseño de endpoints con respuesta $<$5~s. Procesamiento asíncrono y delegación a jobs batch para operaciones largas. \\
Latencia de Supabase Realtime bajo carga alta &
  Retrasos en actualización del tablero del Chef &
  Supabase Realtime plan Pro soporta hasta 500 conexiones simultáneas. Escalable bajo demanda \cite{supabase_pricing}. \\ \rowcolor{altrow}
Acceso a internet requerido en todo momento &
  Sin conectividad, la app no funciona &
  Restricción asumida y documentada. La operación requiere conectividad estable
  del restaurante o datos móviles disponibles en los dispositivos de los usuarios. \\
\hline
\end{tabular}
\end{table}

\noindent\textbf{Veredicto:} El stack tecnológico seleccionado es maduro,
ampliamente adoptado en la industria y con soporte activo. Todas las restricciones
técnicas identificadas tienen soluciones documentadas y probadas en producción
por proyectos de escala comparable. La plataforma puede construirse con el
conocimiento técnico disponible en el mercado mexicano. \textbf{FACTIBLE}.

% -----------------------------------------------------------
\FloatBarrier
\subsection{Factibilidad Operativa}
% -----------------------------------------------------------

La factibilidad operativa evalúa la capacidad de los usuarios finales para
adoptar y operar el sistema propuesto, considerando el perfil tecnológico de
la población objetivo en México.

\subsubsection{Perfil Tecnológico de los Usuarios Finales}

\noindent\textbf{Comensales: adopción de smartphones y QR en México.}
El usuario comensal interactúa con el sistema únicamente a través del navegador
de su dispositivo móvil. Según la Encuesta Nacional sobre Disponibilidad y Uso de
Tecnologías de la Información en los Hogares (ENDUTIH) 2022 del INEGI, el 88.6\%
de la población mexicana de 25 a 34 años es usuaria de internet, y el 89.3\%
accede a internet principalmente desde un teléfono celular \cite{inegi_endutih_2022}. En el segmento de
35 a 44 años, la penetración de internet es del 78.9\%.

El uso de c\'odigos QR para acceder a men\'us digitales se consolid\'o durante
y despu\'es de la pandemia de COVID-19. De acuerdo con la National Restaurant
Association (2023), el 73\% de los comensales de la Generaci\'on~Z y el 46\%
de los Baby~Boomers muestran inter\'es en usar c\'odigos QR para men\'us de
restaurante \cite{nra_state_2023}. A nivel nacional, la alta penetraci\'on de smartphones en
M\'exico \cite{inegi_endutih_2022} respalda la viabilidad de este canal como punto de entrada al sistema.
Esto implica que la curva de aprendizaje para el flujo de activaci\'on de mesa
virtual es m\'inima.

\noindent\textbf{Personal operativo: capacitación requerida.}
Las interacciones del sistema han sido diseñadas para minimizar la curva de
aprendizaje: el mesero solo necesita inicializar mesas y entregar un QR; el chef
solo visualiza un tablero de órdenes; el gerente administra un panel web estándar.
La Tabla~\ref{tab:capacitacion} detalla la curva de aprendizaje estimada por rol.

\begin{table}[H]
\caption{Capacitación Requerida por Rol}
\label{tab:capacitacion}
\centering
\small
\setlength{\tabcolsep}{3pt}
\begin{tabular}{>{\raggedright\arraybackslash}p{2.5cm}
                >{\raggedright\arraybackslash}p{4.5cm}
                >{\raggedright\arraybackslash}p{2.5cm}
                >{\raggedright\arraybackslash}p{6.5cm}}
\hline
\thead{Rol} & \thead{Flujos principales} & \thead{Curva de aprendizaje} & \thead{Capacitación requerida} \\
\hline
Mesero & Inicializar mesa, entregar QR, visualizar estatus, liberar mesa. & 2--4 horas & Sesión práctica de 1~hora + manual de 1~página. \\ \rowcolor{altrow}
Chef & Visualizar tablero de órdenes, actualizar estatus de orden. & 1--2 horas & Sesión práctica de 30~minutos. Interfaz de una sola pantalla. \\
Gerente de Sucursal & Todo lo anterior + gestión de menú, usuarios y dashboard. & 4--8 horas & Capacitación de 2~horas + documentación de usuario. \\ \rowcolor{altrow}
Gerente de Zona / Admin & Gestión multi-sucursal, zonas, analytics. & 6--10 horas & Capacitación de 3~horas + acceso a modo sandbox. \\
Comensal / Anfitrión & Escanear QR, crear contraseña, ordenar, registrar aportación / liquidación interna. & Menos de 5 minutos & Flujo guiado en pantalla y apoyo del mesero en caso necesario. \\
\hline
\end{tabular}
\end{table}

\subsubsection{Análisis de Cambio Organizacional}

La Tabla~\ref{tab:cambio_org} evalúa los factores de cambio que implica la
implantación del sistema y los mecanismos de mitigación propuestos.

\begin{table}[H]
\caption{Análisis de Cambio Organizacional}
\label{tab:cambio_org}
\centering
\small
\setlength{\tabcolsep}{3pt}
\begin{tabular}{>{\raggedright\arraybackslash}p{3.5cm}
                >{\raggedright\arraybackslash}p{5.5cm}
                >{\raggedright\arraybackslash}p{7cm}}
\hline
\thead{Factor de cambio} & \thead{Impacto para el personal} & \thead{Mitigación propuesta} \\
\hline
El comensal genera su propia orden &
  El mesero deja de tomar órdenes en papel. Puede percibirse como reducción de su rol. &
  Comunicar que el mesero se libera para atención de mayor valor. No implica reducción de plantilla. \\ \rowcolor{altrow}
Nuevo flujo de activación de mesa &
  Paso adicional al proceso actual (inicializar mesa virtual antes de sentar comensales). &
  El flujo toma menos de 60~segundos. Onboarding interactivo integrado en la app. \\
Chef visualiza tablero digital &
  Cambio de medio: papel $\rightarrow$ pantalla. Puede generar resistencia en personal de mayor edad. &
  Interfaz con fuente grande. Notificación sonora de nuevas órdenes. Período de transición paralela recomendado. \\ \rowcolor{altrow}
Comensales realizan liquidaciones internas / aportaciones &
  Comensales que prefieren efectivo o pago en caja pueden optar por esa alternativa. &
  El sistema permite registrar aportaciones manuales en caja; la versión 1.0 no procesa pagos bancarios reales. \\
\hline
\end{tabular}
\end{table}

\subsubsection{Contexto de la Industria Restaurantera en México}

Según la Cámara Nacional de la Industria de Restaurantes y Alimentos Condimentados
(CANIRAC), el sector restaurantero en M\'exico concentra m\'as de 700{,}000 unidades
econ\'omicas seg\'un datos del DENUE del INEGI \cite{inegi_denue_2024}. El sector ha experimentado una
aceleración en la adopción de tecnología post-pandemia, especialmente en herramientas de
punto de venta, pedidos digitales, menús QR y presencia en línea \cite{gonzalez_canirac_2024}.

\noindent\textbf{Veredicto:} El perfil tecnológico de los usuarios finales es
compatible con los requerimientos del sistema. El 89.3\% de los mexicanos accede
a internet desde smartphone \cite{inegi_endutih_2022}, y el 64\% ha usado un QR para acceder a un menú
\cite{nra_state_2023}. Los flujos del sistema están diseñados para minimizar la curva
de aprendizaje. La resistencia al cambio en el personal operativo es manejable
con capacitación mínima y transición gradual. \textbf{FACTIBLE}.

% -----------------------------------------------------------
\FloatBarrier
\subsection{Factibilidad Económica}

\subsubsection{Criterio de Estimaci\'on Econ\'omica}
Para evitar mezclar gasto efectivo con valor de trabajo desarrollado, la factibilidad
econ\'omica distingue dos conceptos:
\begin{itemize}
  \item \textbf{Costo operativo real:} pagos necesarios para mantener la infraestructura del sistema (Vercel, Supabase, VPS, dominio), que constituyen el gasto recurrente real.
  \item \textbf{Valor equivalente de desarrollo:} estimaci\'on del costo que implicar\'ia contratar externamente el desarrollo; se incluye como referencia de esfuerzo y no como presupuesto efectivo de la tesis.
\end{itemize}

\subsubsection{Metodología de Estimación de Costos}

El presente estudio de factibilidad econ\'omica aplica la metodolog\'ia de
\textbf{Estimaci\'on Bottom-Up}, definida por el Project Management Institute
como la descomposici\'on del proyecto en actividades individuales con estimaci\'on
de costo independiente para cada una, que posteriormente se suman para obtener
el presupuesto total \cite{stps_observatorio}. Este m\'etodo es id\'oneo para el presente proyecto
porque los m\'odulos (OE-01 a OE-13) est\'an claramente delimitados y permiten
estimar horas de desarrollo por objetivo espec\'ifico con alta precisi\'on.

Los par\'ametros base utilizados son: tarifa de \$500~MXN/hora para un desarrollador
\textit{full-stack senior} con el perfil TypeScript~+~Next.js~+~Supabase, consistente
con el rango medio del mercado mexicano de desarrollo de software en 2026 seg\'un
el Observatorio Laboral de la STPS \cite{stps_observatorio}; y un tipo de cambio referencial de
\$17.50~MXN/USD conforme a Banxico. Ambos par\'ametros son editables en el archivo
de gesti\'on de costos adjunto (\textit{Gesti\'on\_Costos.xlsx}).

\subsubsection{Estimación de Costos de Desarrollo (Inversión Inicial)}

La Tabla~\ref{tab:costos_dev} resume el valor econ\'omico equivalente de desarrollo
por m\'odulo, derivado del desglose bottom-up de los objetivos espec\'ificos y las
actividades de la Fase~1 (an\'alisis y dise\~no). El detalle completo, con
horas por actividad y f\'ormulas editables, se encuentra en la hoja
\textit{2.~Estimaci\'on Bottom-Up} del archivo Excel adjunto.

\setstretch{1.2}
\setlength{\tabcolsep}{4pt}
\begin{longtable}{>{\raggedright\arraybackslash}p{3.5cm}
                >{\raggedright\arraybackslash}p{5.5cm}
                >{\centering\arraybackslash}p{1.6cm}
                >{\centering\arraybackslash}p{2.5cm}
                >{\centering\arraybackslash}p{2cm}}
\caption{Estimación de Costos de Desarrollo por Módulo (Bottom-Up)} \label{tab:costos_dev} \\
\hline
\thead{Módulo / Fase} & \thead{Actividades principales} &
\thead{Horas est.} & \thead{Costo RRHH\,(MXN)} & \thead{Fase} \\
\hline
\endfirsthead
\caption[]{TABLA \thetable{} (Continuación)} \\
\hline
\thead{Módulo / Fase} & \thead{Actividades principales} &
\thead{Horas est.} & \thead{Costo RRHH\,(MXN)} & \thead{Fase} \\
\hline
\endhead
\hline
\endfoot
OE-01 Infraestructura & Vercel CI/CD, schema Prisma, Supabase base & 160 & \$80,000 & 2 \\ \rowcolor{altrow}
OE-02 Estructura Org. & CRUD Cadena/Zona/Sucursal + RLS & 60 & \$30,000 & 2 \\
OE-03 Auth \& Roles & Supabase Auth + RLS policies + control de sesiones & 100 & \$50,000 & 2 \\ \rowcolor{altrow}
OE-04 Gestión Usuarios & Alta usuarios + contrase\~na temporal & 50 & \$25,000 & 2 \\
OE-05 Cat\'alogo & CRUD categor\'ias + productos + variantes & 90 & \$45,000 & 2 \\ \rowcolor{altrow}
OE-06/07 Mesas & Mesas f\'isicas + mesa virtual + QR HMAC & 150 & \$75,000 & 2 \\
OE-08/09 \'Ordenes + Chef & Canasta, \'ordenes, tablero KDS Realtime & 170 & \$85,000 & 2 \\ \rowcolor{altrow}
OE-10/11 Liquidaci\'on & Concurrencia \texttt{SELECT FOR UPDATE} + constancia interna & 140 & \$70,000 & 2 \\
OE-12 Dashboard & Pipeline Spark (Medallion) + Dashboard multi-nivel & 120 & \$60,000 & 2 \\ \rowcolor{altrow}
OE-13 Seguridad & Rate limiting + RLS audit + pruebas carga & 80 & \$40,000 & 2 \\
Pruebas QA & Unitarias + integraci\'on + carga (k6) & 80 & \$40,000 & 2 \\ \rowcolor{altrow}
Fase 1 --- An\'alisis y Dise\~no (Natalia) & Marco te\'orico, flujos, UML, BD, wireframes, TT & 660 & \$330,000 & 1 \\
Gesti\'on de Proyecto (ambos) & Ceremonias Scrum, documentaci\'on, retrospectivas & 80 & \$40,000 & 1/2 \\
\hline
\rowcolor{green}
\textbf{TOTAL RRHH (valor equivalente)} & & \textbf{1,940} & \textbf{\$970,000} & --- \\
\end{longtable}
\setstretch{1.5}

\subsubsection{Costos de Infraestructura Cloud (Operación Mensual)}

Los costos de infraestructura son recurrentes y se calculan con base en las tarifas
p\'ublicas de cada servicio. La Tabla~\ref{tab:costos_infra} presenta los costos
mensuales y anuales para el primer a\~no del proyecto (12 meses operativos).

\begin{table}[H]
\caption{Costos de Infraestructura Cloud (Año 1 — 12 Meses)}
\label{tab:costos_infra}
\centering
\small
\setlength{\tabcolsep}{3pt}
\begin{tabular}{>{\raggedright\arraybackslash}p{3cm}
                >{\raggedright\arraybackslash}p{2.5cm}
                >{\centering\arraybackslash}p{1.5cm}
                >{\centering\arraybackslash}p{1.5cm}
                >{\centering\arraybackslash}p{2.5cm}
                >{\raggedright\arraybackslash}p{5cm}}
\hline
\thead{Servicio} & \thead{Plan} & \thead{USD/mes} &
\thead{MXN/mes} & \thead{MXN\,(12 meses)} & \thead{Referencia} \\
\hline
Vercel & Pro (1 seat) & \$20 & \$350 & \$4,200 & Vercel Pricing \cite{vercel_pricing} \\ \rowcolor{altrow}
Supabase & Pro & \$25 & \$437.50 & \$5,250 & Supabase Pricing \cite{supabase_pricing} \\
VPS Hostinger KVM 2 & 2 vCPU, 8 GB RAM, 100 GB NVMe, 8 TB BW & \$8.99 & \$157.33 & \$1,887.96 & Hostinger KVM 2 \cite{hostinger_kvm2} \\ \rowcolor{altrow}
Dominio .com.mx & Anual & \$0.83 & \$14.53 & \$174.36 & NIC M\'exico / GoDaddy \\
\hline
\rowcolor{green}
\textbf{Total infraestructura (promocional)} & & \textbf{\$54.82} & \textbf{\$959.36} & \textbf{\$11,512.32} & \\
\hline
\rowcolor{altrow}
VPS Hostinger KVM 2 (renovaci\'on) & Renovaci\'on & \$14.99 & \$262.33 & \$3,147.96 & Precio de renovaci\'on indicado en proveedor \\
\hline
\rowcolor{green}
\textbf{Total infraestructura (renovaci\'on)} & & \textbf{\$60.82} & \textbf{\$1,064.36} & \textbf{\$12,772.32} & \\
\hline
\end{tabular}
\end{table}

\subsubsection{Valor Económico Equivalente del Proyecto}

La Tabla~\ref{tab:presupuesto_total} presenta el valor econ\'omico equivalente del
proyecto durante el primer a\~no. Esta estimaci\'on no representa un gasto efectivo
del Trabajo Terminal, sino el costo aproximado que tendr\'ia desarrollar una soluci\'on
similar si se contrataran recursos humanos, consultor\'ia y herramientas en condiciones
de mercado. La reserva de contingencia del 15\% se conserva como referencia de planeaci\'on
\cite{pmi_pmbok_7}.

\begin{table}[H]
\caption{Valor Económico Equivalente del Proyecto — Año 1}
\label{tab:presupuesto_total}
\centering
\small
\setlength{\tabcolsep}{4pt}
\begin{tabular}{>{\raggedright\arraybackslash}p{7cm}
                >{\centering\arraybackslash}p{2cm}
                >{\centering\arraybackslash}p{3.5cm}
                >{\centering\arraybackslash}p{3.5cm}}
\hline
\thead{Concepto} & \thead{Tipo} & \thead{Costo MXN} & \thead{Costo USD} \\
\hline
Recursos Humanos --- Desarrollo (Diego) & Directo/Fijo & \$620,000 & \$35,429 \\ \rowcolor{altrow}
Recursos Humanos --- An\'alisis y Dise\~no (Natalia) & Directo/Fijo & \$330,000 & \$18,857 \\
Gesti\'on de Proyecto (ambos) & Directo/Fijo & \$40,000 & \$2,286 \\ \rowcolor{altrow}
Infraestructura cloud (12 meses) & Directo/Fijo & \$11,512 & \$658 \\
Capacitaci\'on del personal (onboarding) & Directo/Fijo & \$12,000 & \$686 \\ \rowcolor{altrow}
Consultor\'ia legal (LFPDPPP + T\'erminos) & Indirecto/Fijo & \$25,000 & \$1,429 \\
Herramientas de dise\~no (Figma Pro 6 meses) & Directo/Variable & \$3,150 & \$180 \\ \rowcolor{altrow}
Herramientas de documentaci\'on y PM & Indirecto/Variable & \$1,750 & \$100 \\
\hline
\rowcolor{orange}
\textbf{Subtotal Base} & & \textbf{\$1,043,412} & \textbf{\$59,623} \\ \rowcolor{orange}
Reserva de Contingencia (15\% --- PMI PMBOK) & Buffer & \$156,512 & \$8,944 \\
\hline
\rowcolor{green}
\textbf{VALOR EQUIVALENTE TOTAL A\~NO 1} & & \textbf{\$1,199,924 MXN} & \textbf{\$68,567 USD} \\
\hline
\end{tabular}
\end{table}

\subsubsection{Comparación Referencial frente a Alternativas SaaS}

La comparaci\'on no busca demostrar sustituci\'on funcional completa frente a plataformas
POS comerciales, ya que dichas soluciones incluyen m\'odulos fuera del alcance de la
versi\'on 1.0, como inventario avanzado, facturaci\'on, pagos bancarios, soporte t\'ecnico
comercial y hardware especializado. Su prop\'osito es ubicar el costo operativo recurrente
del prototipo frente a servicios SaaS relacionados.

\begin{table}[H]
\caption{Comparación Referencial frente a Alternativas SaaS}
\label{tab:build_buy}
\centering
\small
\setlength{\tabcolsep}{3pt}
\begin{tabular}{>{\raggedright\arraybackslash}p{3.5cm}
                >{\centering\arraybackslash}p{2.8cm}
                >{\centering\arraybackslash}p{3.2cm}
                >{\centering\arraybackslash}p{2cm}
                >{\raggedright\arraybackslash}p{4.2cm}}
\hline
\thead{Solución} & \thead{Costo mensual\,(MXN)} & \thead{Costo anual\,(MXN)} &
\thead{Propie-dad código} & \thead{Observación} \\
\hline
\textbf{El presente sistema} & \textbf{\$959} & \textbf{\$11,512} & \textbf{Sí} &
  Solo infraestructura recurrente; el desarrollo es valor equivalente no recurrente. \\ \rowcolor{altrow}
Toast POS + Mobile Order~\&~Pay & \$2,500 & \$90,000 & No &
  Sin mesa virtual con contrase\~na. \\
Square for Restaurants + KDS & \$2,200 & \$79,200 & No &
  Split bill desde cajero, no desde comensal. \\ \rowcolor{altrow}
Lightspeed Restaurant & \$3,500 & \$126,000 & No &
  Sin jerarqu\'ia cadena/zona/sucursal. \\
Yumminn & \$2,000 & \$72,000 & No &
  Sin KDS propio ni cuentas individuales. \\ \rowcolor{altrow}
Oracle MICROS POS & \$15,000 & \$540,000 & No &
  Enterprise. Hardware adicional requerido. \\
\hline
\end{tabular}
\end{table}

\noindent\footnotesize{\textit{Nota:} Los costos de soluciones comerciales son aproximados
y pueden variar según país, número de sucursales, módulos contratados, hardware,
soporte, comisiones e integraciones. La comparación es referencial y no implica
equivalencia funcional completa.}
\normalsize

\noindent El costo de infraestructura mensual del sistema es sensiblemente inferior al
costo operativo típico de soluciones SaaS comerciales. Sin embargo, la estimaci\'on del
"valor equivalente de desarrollo" (horas multiplicadas por tarifas de mercado) se presenta
solamente como referencia para dimensionar el esfuerzo t\'ecnico y no representa gasto
efectivo en el contexto de un Trabajo Terminal.


\noindent\textbf{Veredicto:} El proyecto es factible para un contexto acad\'emico y de
validaci\'on t\'ecnica. Manteniendo la exclusi\'on de procesamiento bancario real en la
versi\'on 1.0, los costos operativos son bajos y la mayor inversi\'on corresponde al
esfuerzo de desarrollo (valor equivalente). \textbf{FACTIBLE}.

% -----------------------------------------------------------
\FloatBarrier
\subsection{Factibilidad Legal}
% -----------------------------------------------------------

La factibilidad legal evalúa el cumplimiento del marco normativo mexicano e
internacional aplicable a la plataforma, incluyendo protección de datos personales,
registro interno de liquidaciones, facturación electrónica como trabajo futuro,
licencias de software y normativas de telecomunicaciones.

\subsubsection{Ley Federal de Protección de Datos Personales en Posesión de
Particulares (LFPDPPP)}

La plataforma recolecta y procesa datos personales de los usuarios colaboradores
(correo electrónico, nombre, historial de accesos) a través de Supabase Auth.
Adicionalmente, se registran datos de comportamiento de comensales (órdenes
generadas, montos registrados/aportaciones) aunque de forma no nominal si el comensal no se
registra formalmente. Este tratamiento está sujeto a la LFPDPPP, publicada en
el Diario Oficial de la Federación (DOF) el 5 de julio de 2010, y a su Reglamento,
publicado el 21 de diciembre de 2011 \cite{lfpdppp}.

Las sanciones por incumplimiento van desde 100 hasta 320,000 días de salario
mínimo general vigente (SMGV), equivalentes a aproximadamente \$12,200 hasta
\$39,040,000 MXN con el SMGV 2024 de \$248.93 MXN/día \cite{reglamento_lfpdppp}.

La Tabla~\ref{tab:lfpdppp} detalla las obligaciones aplicables y su estado de
cumplimiento en la versión 1.0 de la plataforma.

\setstretch{1.2}
\setlength{\tabcolsep}{4pt}
\begin{longtable}{>{\raggedright\arraybackslash}p{3.5cm}
                >{\centering\arraybackslash}p{1.5cm}
                >{\raggedright\arraybackslash}p{6.5cm}
                >{\raggedright\arraybackslash}p{4cm}}
\caption{Obligaciones LFPDPPP Aplicables a la Plataforma} \label{tab:lfpdppp} \\
\hline
\thead{Obligación legal} & \thead{Artículo} & \thead{Acción requerida} & \thead{Estado en v1.0} \\
\hline
\endfirsthead
\caption[]{TABLA \thetable{} (Continuación)} \\
\hline
\thead{Obligación legal} & \thead{Artículo} & \thead{Acción requerida} & \thead{Estado en v1.0} \\
\hline
\endhead
\hline
\endfoot
Aviso de Privacidad &
  Art.\ 15--16 &
  Redactar y publicar el Aviso de Privacidad antes del primer tratamiento de datos. Incluir identidad del responsable, finalidades, derechos ARCO y transferencias a Supabase/Vercel. &
  \textbf{Requerido} antes del lanzamiento. \\ \rowcolor{altrow}
Consentimiento del titular &
  Art.\ 8 &
  Obtener consentimiento expreso al crear cuenta o acceder por primera vez. &
  Implementar checkbox de aceptación en el registro. \\
Principio de licitud &
  Art.\ 7 &
  Los datos solo pueden tratarse para las finalidades declaradas en el Aviso de Privacidad. &
  Cumplido por diseño: datos usados solo para operar el servicio. \\ \rowcolor{altrow}
Principio de minimización &
  Art.\ 7 &
  Solo recolectar los datos estrictamente necesarios. &
  Cumplido: solo correo y nombre de empleados. Comensales son anónimos. \\
Derechos ARCO &
  Art.\ 28--37 &
  Implementar mecanismo para acceso, rectificación, cancelación y oposición. Plazo: 20 días hábiles. &
  Módulo de eliminación de cuenta requerido. \\ \rowcolor{altrow}
Medidas de seguridad &
  Art.\ 19 &
  TLS, cifrado en tránsito y en reposo, RLS en PostgreSQL, acceso mínimo por rol. &
  Implementado en la arquitectura del sistema. \\
Transferencias internacionales &
  Art.\ 36 &
  Datos alojados en Supabase (AWS us-east) están fuera de México. Debe declararse en el Aviso. &
  Requiere cláusula específica en el Aviso de Privacidad. \\ \rowcolor{altrow}
Responsable de datos &
  Art.\ 3 fracc.\ XIV &
  Designar un responsable de datos (Data Controller) dentro de la organización. &
  Designar antes del lanzamiento. \\
\end{longtable}
\setstretch{1.5}

\subsubsection{PCI DSS y Exclusión de Pagos Bancarios en v1.0}

La versi\'on~1.0 de la plataforma no almacena, transmite ni procesa datos de tarjeta.
La liquidaci\'on se maneja únicamente como un registro interno del estado de cuenta dentro de la
aplicación. Por tanto, el alcance PCI DSS queda fuera del sistema actual.

La integraci\'on con agregadores de pago se documenta como trabajo futuro. En caso de
incorporarse, deberá realizarse mediante un proveedor certificado que permita reducir el
alcance PCI de la plataforma, evitando que el sistema vea, transmita o almacene datos de
tarjeta.

\subsubsection{Obligaciones Fiscales: CFDI}

El Artículo 29 del Código Fiscal de la Federación (CFF) establece la obligación
de expedir Comprobantes Fiscales Digitales por Internet (CFDI) por los actos o
actividades de los contribuyentes \cite{cff_art29}. La plataforma en su versión 1.0 no emite
CFDI directamente; la obligación recae sobre la cadena restaurantera. No obstante,
se recomienda evaluar la integración con un Proveedor Autorizado de Certificación
(PAC) del SAT \cite{sat_pac} para automatizar la emisión de CFDI a solicitud del comensal,
especialmente en el segmento empresarial donde el 78\% de los comensales
de restaurantes de servicio completo requieren factura \cite{inegi_canirac_2023}.

\subsubsection{Ley Federal de Protección al Consumidor (LFPC)}

La plataforma participa en el registro digital de órdenes y liquidaciones internas
asociadas a una relación de consumo en restaurante, por lo que se consideran
principios de transparencia, información clara, términos y condiciones, y mecanismos
de aclaración conforme a la LFPC \cite{lfpc}.

La Tabla~\ref{tab:lfpc} detalla las obligaciones aplicables y su implementación
en la plataforma.

\begin{table}[H]
\caption{Obligaciones LFPC Aplicables a la Plataforma}
\label{tab:lfpc}
\centering
\small
\setlength{\tabcolsep}{3pt}
\begin{tabular}{>{\raggedright\arraybackslash}p{4cm}
                >{\centering\arraybackslash}p{2cm}
                >{\raggedright\arraybackslash}p{10cm}}
\hline
\thead{Obligación LFPC} & \thead{Artículo} & \thead{Implementación en la plataforma} \\
\hline
Información clara sobre consumo y liquidación &
  Art.\ 7 &
  La pantalla de liquidación muestra el desglose completo del consumo y el monto pendiente antes de confirmar cualquier registro interno. \\ \rowcolor{altrow}
Confirmación explícita del pedido &
  Art.\ 76 BIS fracc.\ II &
  Diálogo de confirmación antes de enviar cada orden. Confirmación final antes de registrar la liquidación interna. \\
Derecho a cancelar antes de confirmar &
  Art.\ 76 BIS fracc.\ III &
  El comensal puede vaciar su canasta o cancelar el registro interno antes de confirmarlo. \\ \rowcolor{altrow}
Términos y condiciones accesibles &
  Art.\ 76 BIS fracc.\ I &
  Términos de Uso publicados y aceptados al acceder por primera vez a la plataforma. \\
Mecanismo de quejas &
  Art.\ 24 &
  Publicar correo o formulario de atención al cliente. PROFECO puede recibir quejas sobre el servicio. \\
\hline
\end{tabular}
\end{table}

\subsubsection{Licencias de Software del Stack}

El análisis de licencias garantiza que el uso comercial de cada componente del
stack sea legal. La Tabla~\ref{tab:licencias} muestra que ningún componente impone
restricciones incompatibles con el modelo de negocio SaaS de la plataforma \cite{osi_licenses}.

\begin{table}[H]
\caption{Análisis de Licencias del Stack Tecnológico}
\label{tab:licencias}
\centering
\small
\setlength{\tabcolsep}{3pt}
\begin{tabular}{>{\raggedright\arraybackslash}p{3cm}
                >{\raggedright\arraybackslash}p{3cm}
                >{\centering\arraybackslash}p{2.5cm}
                >{\centering\arraybackslash}p{2.5cm}
                >{\raggedright\arraybackslash}p{5cm}}
\hline
\thead{Componente} & \thead{Licencia} & \thead{¿Uso comercial?} &
\thead{¿Publicar código?} & \thead{Restricciones relevantes} \\
\hline
Next.js & MIT License & Sí & No & Mantener aviso de copyright. \\ \rowcolor{altrow}
React & MIT License & Sí & No & Mantener aviso de copyright. \\
TypeScript & Apache 2.0 & Sí & No & Incluir copia de la licencia en distribuciones. \\ \rowcolor{altrow}
Prisma ORM & Apache 2.0 & Sí & No & Incluir copia de la licencia en distribuciones. \\
Supabase (cliente JS) & MIT License & Sí & No & Mantener aviso de copyright. \\ \rowcolor{altrow}
PostgreSQL & PostgreSQL License & Sí & No & Sin restricciones comerciales. \\
Node.js & MIT License & Sí & No & Mantener aviso de copyright. \\ \rowcolor{altrow}
Vercel (plataforma) & Propietario (SaaS) & Sí (Plan Pro) & N/A & Sujeto a Términos de Servicio de Vercel. \\
Supabase (plataforma) & Propietario + AGPLv3 & Sí (Plan Pro cloud) & Solo en self-hosting & En modo BaaS cloud, el código cliente no está sujeto a AGPL \cite{supabase_license}. \\
\hline
\end{tabular}
\end{table}

\subsubsection{Normativas de Telecomunicaciones}

La plataforma es un servicio de software entregado a través de internet (SaaS/PWA)
y está regulada por la Ley Federal de Telecomunicaciones y Radiodifusión (LFTR) \cite{lftr}.
La plataforma no opera como concesionaria de red pública ni como proveedor de
servicios de telecomunicaciones (voz, SMS), por lo que la LFTR no impone
obligaciones de registro o concesión. No obstante, si en versiones futuras se
implementan notificaciones push vía SMS, se deberá revisar el Art. 191 de la
LFTR relativo a la identidad del remitente en comunicaciones comerciales.

\noindent\textbf{Veredicto:} El proyecto es legalmente viable. Las acciones legales
no negociables antes del lanzamiento son: (1) redacci\'on y publicaci\'on del Aviso de
Privacidad (Arts.~15--16 LFPDPPP); (2) implementaci\'on de consentimiento expreso y
mecanismo ARCO; y (3) T\'erminos y Condiciones conforme al Art.~76~BIS de la LFPC.
La decisi\'on de no integrar un agregador de pago v\'ia tarjeta en v1.0 evita el riesgo
RL-02 (PCI~DSS) dentro del alcance de la versión 1.0. El stack utiliza exclusivamente
licencias permisivas (MIT, Apache~2.0, BSD) sin restricciones sobre el uso comercial.
\textbf{FACTIBLE CON ACCIONES REQUERIDAS}.

